\documentclass[12pt,a4paper]{article}
\usepackage[english]{babel}
\usepackage{titlesec}

% \titleformat{\chapter} % command [display] % shape {\bfseries\huge} % format
%   {Chapter \thechapter} % label {0pt} % sep {} % before-code


\usepackage[T5]{fontenc}
\usepackage{mathptmx}[ptm]
\usepackage{a4wide, amssymb, epsfig, latexsym, array, hhline, fancyhdr}
\usepackage[normalem]{ulem}
%\usepackage{soul}
\usepackage{float}
\usepackage[makeroom]{cancel}
\usepackage{amsmath}
\usepackage{amsthm}
\usepackage{multicol, longtable, amscd}
\usepackage{diagbox}%Make diagonal lines in tables
\usepackage{booktabs}
\usepackage{alltt}
\usepackage[framemethod=tikz]{mdframed}% For highlighting paragraph backgrounds
\usepackage{indentfirst}
\usepackage{caption,subcaption}
\usepackage{svg}
\graphicspath{ {figures/} }
\usepackage{lastpage}
\usepackage[lined,boxed,commentsnumbered]{algorithm2e}
\usepackage{enumerate}
\usepackage{color}
\usepackage{graphicx}% Standard graphics package
\graphicspath{ {figures/} }
\usepackage{array}
\usepackage{tabularx, caption}
\usepackage{tabularray}
\usepackage{multirow}
\usepackage{multicol}
\usepackage{rotating}
\usepackage{graphics}
\usepackage{geometry}
\usepackage{setspace}
\usepackage{epsfig}
\usepackage{tikz}
\usepackage[numbers]{natbib}
\usepackage{subcaption}
\usepackage{dirtree}

\usetikzlibrary{arrows, snakes, backgrounds, calc}
\usepackage[unicode]{hyperref}
\hypersetup{urlcolor=blue,linkcolor=black,citecolor=black,colorlinks=true} 

\usepackage{xcolor}

\definecolor{codegreen}{rgb}{0,0.6,0}
\definecolor{codegray}{rgb}{0.5,0.5,0.5}
\definecolor{codepurple}{rgb}{0.58,0,0.82}
\definecolor{backcolour}{rgb}{0.95,0.95,0.92}

\usepackage{listings}
\lstdefinestyle{mystyle}{ backgroundcolor=\color{backcolour},   
    commentstyle=\color{codegreen}, keywordstyle=\color{magenta},
    numberstyle=\tiny\color{codegray}, stringstyle=\color{codepurple},
    basicstyle=\ttfamily\footnotesize, breakatwhitespace=false,         
    breaklines=true,                 
    captionpos=b,                    
    keepspaces=true,                 
    numbers=left,                    
    numbersep=5pt,                  
    showspaces=false,                
    showstringspaces=false, showtabs=false,                  
    tabsize=2 }
\lstset{style=mystyle}



\def\thesislayout{	% A4: 210 × 297
	\geometry{ a4paper, total={160mm,240mm},  % fix over page
		left=30mm, top=30mm, } }
\thesislayout

%\usepackage{fancyhdr}
\setlength{\headheight}{40pt}
\pagestyle{fancy}
\fancyhead{} % clear all header fields
\fancyhead[L]{
 \begin{tabular}{rl}
    \begin{picture}(25,15)(0,0) \put(0,-8){\includegraphics[width=10mm,
    height=10mm]{Images/hcmut.png}}
    %\put(0,-8){\epsfig{width=10mm,figure=hcmut.eps}}
   \end{picture}&
	%\includegraphics[width=8mm, height=8mm]{hcmut.png} & %
	\begin{tabular}{l}
        \textbf{ Ho Chi Minh City University of Technology } \\
        \textbf{ Faculty of Computer Science and Engineering }
	\end{tabular} 	
 \end{tabular}
} \fancyhead[R]{
	\begin{tabular}{l}
		\tiny \bf \\
		\tiny \bf \end{tabular}  } \fancyfoot{} % clear all footer fields
\fancyfoot[L]{\scriptsize Data Engineering - Academic year 2024 - 2025}
\fancyfoot[R]{\scriptsize  Page {\thepage}/\pageref{LastPage}}
\renewcommand{\headrulewidth}{0.3pt}
\renewcommand{\footrulewidth}{0.3pt}

%%%%%%%%%%%%%%%%%%%%%%%%%%%%%%%%%%%%%%%%%%
% Tạo một môi trường mới cho các bảng đặc tả UC, giúp định dạng kiểu
\newenvironment{usecase_table}[1][Test]
{
\begin{table}[H]
% Left-Right Table padding
    \setlength{\tabcolsep}{6pt}
    \renewcommand{\arraystretch}{1.5}
    % \begin{tabularx}{\textwidth}{bs}
    \def\savedcaption{\caption{#1}}%
    \begin{tabular}{|>{\bfseries} p{0.2\linewidth}|p{0.7\linewidth}|}

} {
    \end{tabular}
    \savedcaption
\end{table}
}

\newenvironment{usecase_enum}
{
    \begin{enumerate*}[itemjoin={\newline}]
} {
    \end{enumerate*}
}

\makeatletter
\newenvironment{acknowledgments}{
	\small
	\begin{center}
	  {\bfseries Acknowledgments\vspace{-.5em}\vspace{\z@}}
	\end{center}
	\quotation
  }{}
\makeatother

% Declaration
\makeatletter
\newenvironment{declaration}{
	\small
	\begin{center}
	  {\bfseries Declaration\vspace{-.5em}\vspace{\z@}}
	\end{center}
	\quotation
  }{}
\makeatother

\makeatletter
\newenvironment{abstr}{
        \small
        \begin{center}
	{\bfseries Abstract \vspace{-.5em}\vspace{\z@}}
        \end{center}
        \quotation
    }{}
\makeatother

\makeatletter
\newenvironment{abstrsss}{
        \small
        \begin{center}
	{\bfseries Abstractss \vspace{-.5em}\vspace{\z@}}
        \end{center}
        \quotation
    }{}
\makeatother
%%%%%%%%%%%%%%%%%%%%%%%%%%%%%%%%%%%%%%%%%%%%%%5
\setcounter{secnumdepth}{4}
\setcounter{tocdepth}{3}
\makeatletter
\newcounter {subsubsubsection}[subsubsection]
\renewcommand\thesubsubsubsection{\thesubsubsection .\@alph\c@subsubsubsection}
\newcommand\subsubsubsection{\@startsection{subsubsubsection}{4}{\z@}%
                                     {-3.25ex\@plus -1ex \@minus -.2ex}%
                                     {1.5ex \@plus .2ex}%
                                     {\normalfont\normalsize\bfseries}}
\newcommand*\l@subsubsubsection{\@dottedtocline{3}{10.0em}{4.1em}}
\newcommand*{\subsubsubsectionmark}[1]{}
\makeatother



\everymath{\color{black}}%make in-line maths symbols blue to read/check easily

\sloppy
\captionsetup[figure]{labelfont={small,bf},textfont={small,it},belowskip=-1pt,aboveskip=-9pt}
%space removed between caption, figure, and text
\captionsetup[table]{labelfont={small,bf},textfont={small,it},belowskip=-1pt,aboveskip=7pt}
\setlength{\floatsep}{5pt plus 2pt minus 2pt}
\setlength{\textfloatsep}{5pt plus 2pt minus 2pt}
\setlength{\intextsep}{10pt plus 2pt minus 2pt}


% Declare a variable named Proc using \newcommand

\newcommand{\Uni}{Ho Chi Minh University - Vietnam National Unversity, Ho Chi Minh City}

\renewcommand{\labelitemi}{\normalfont\bfseries\textendash}
\renewcommand{\labelitemii}{\textbullet}
\renewcommand*\baselinestretch{1.5}\selectfont


\thesislayout

\begin{document}

\begin{titlepage}
    \begin{tikzpicture}[remember picture, overlay]
        \draw[line width = 4pt] ($(current page.north west) + (1.0in,-0.5in)$)
        rectangle ($(current page.south east) + (-0.4in,0.5in)$);
        \draw[line width=1.5pt]
        ($ (current page.north west) + (1.05in,-0.55in) $) rectangle ($ (current
        page.south east) + (-0.45in,0.55in) $);
    \end{tikzpicture}

    \begin{center}
        \large \textbf{VIETNAM NATIONAL UNIVERSITY, HO CHI MINH CITY} \\
        \large \textbf{HO CHI MINH UNIVERSITY OF TECHNOLOGY} \\
        \large \textbf{FACULTY OF COMPUTER SCIENCE AND ENGINEERING}
    \end{center}

    \begin{figure}[h!]
        \begin{center}
            \includegraphics[width=5cm]{Images/hcmut.png}
        \end{center}
    \end{figure}

    \begin{center}
        \begin{tabular}{c}
        \multicolumn{1}{l}{\textbf{{\Large Data Engineering}}}\\
        ~~\\
        \hline
        \\
        \multicolumn{1}{l}{\textbf{{\Large Assignment Report}}}\\
        \\
        \textbf{{\Huge Flight Operation Analysis}}\\
        \\
        \hline
        \end{tabular}
        \end{center}

    \vspace{1cm}
    \begin{table}[H]
        \begin{tabular}{rrl}
        \hspace{5 cm} & \textbf{Mentor}: & Võ Thị Ngọc Châu\\
        
        & \textbf{Student}: & Trần Hà Tuấn Kiệt -- 2011493 \\
        & & Lã Minh Đức -- 2110132 \\
        & & Cao Đức Dương -- 2110971 \\
        & & Nguyễn Minh Thuấn -- 2010663 \\
        
        \end{tabular}
        \end{table}
    \vspace{1cm}

    \begin{center}
        {\large Ho Chi Minh City, \the\month/\the\year}
    \end{center}
\end{titlepage}


\newpage

%\thispagestyle{empty} \chapter*{Work Assignment} \begin{table}[H] \centering
% \renewcommand{\arraystretch}{1.5} \begin{tabular}{|c|c|c|l|} \hline
% \textbf{No.}       & \textbf{Full name}                 & \textbf{Student ID}
% & \textbf{Work Assignment}                     \\
%         \hline %%%%%Student 1%%%%%%%%%%% \multirow{4}{*}{1} &
%         \multirow{4}{*}{Trần Hà Tuấn Kiệt} & \multirow{4}{*}{2011493} & -
%         Exploratory Data Analysis                  \\
%                            &                                    &
%                            & - System proposal and architecture           \\
%                            &                                    &
%                            & - Research and Develop with new technologies \\
%                            &                                    &
%                            & - Implement system endpoints                 \\
%         \hline %%%%%Student 2%%%%%%%%%%% \multirow{4}{*}{2} &
%         \multirow{4}{*}{Nguyễn Đức Thuỵ}   & \multirow{4}{*}{2012158} & -
%         Exploratory Data Analysis                  \\
%                            &                                    &
%                            & - System proposal and architecture           \\
%                            &                                    &
%                            & - Research and Develop with new technologies \\
%                            &                                    &
%                            & - Implement data processing flows            \\
%         \hline \end{tabular} \end{table} \newpage

% Revert title name back to English
\renewcommand{\contentsname}{Table of Contents}
\renewcommand{\listfigurename}{List of Figures}
\renewcommand{\listtablename}{List of Tables}
\renewcommand{\bibname}{References}
\renewcommand{\figurename}{Figure}

\tableofcontents
\newpage

%\thispagestyle{empty}
% \renewcommand\thechapter{Chapter \arabic{chapter}}

\section{Introduction about Aviation Industry}
The aviation industry is a critical sector that relies heavily on data and
technology to maintain high levels of operational efficiency, ensure optimal
fleet and crew scheduling, and support data-driven decision-making. With some
modern aircraft capable of generating up to 20 terabytes of data per hour, the
industry faces both opportunities and challenges in handling vast amounts of
information. This data supports crucial operations managed by organizations like
the FAA's Air Traffic Organization, which coordinates services for over 45,000
flights daily. In this presentation, we will explore the impact of data on
aviation efficiency and highlighting the industry's reliance on advanced data
management and analysis.

\section{Problem statement}
\subsection{Description}
The aviation industry is currently facing a set of interconnected challenges,
limitations, and impacts that hinder its ability to operate efficiently and
maintain customer satisfaction. These issues stem primarily from data management
difficulties, limited predictive capabilities, and the resulting operational
inefficiencies.
    \begin{itemize}
        \item \textbf{Current Challenges:} The industry struggles with
        inefficient data handling and analysis, which leads to delays in
        decision-making. With modern aircraft generating vast amounts of data,
        the ability to process and interpret this information in real-time is
        critical. However, existing systems often fail to keep up, resulting in
        bottlenecks and slower response times.
        \item \textbf{Limitations:} One of the main constraints in the aviation
        sector is the inability to process large volumes of data in real-time.
        This limitation restricts the use of predictive analytics, which could
        otherwise help in foreseeing and mitigating potential operational
        issues. The lack of predictive analytics means that the industry remains
        largely reactive rather than proactive in managing its operations.
        \item \textbf{Impact:} The combined effect of these challenges and
        limitations leads to increased operational costs, as inefficiencies
        drive up expenses. Additionally, customer satisfaction is negatively
        affected due to delays and a lack of seamless service. When data
        processing and decision-making are delayed, passengers experience
        disruptions, which can diminish trust and satisfaction with airline
        services.
    
    \end{itemize}
    $\quad$This problem statement underscores the need for innovative solutions
    that can handle and analyze large data volumes efficiently, enabling the
    industry to transition to predictive, data-driven decision-making for better
    operational outcomes and enhanced customer experience.
    \subsection{Objectives}
    To address the challenges and limitations in the aviation industry, several
    key objectives have been identified to enhance efficiency, reduce costs, and
    harness the full potential of data:
    \begin{itemize}
        \item \textbf{Real-time Data Processing:} Achieving real-time data
        processing capabilities is essential for the aviation industry to make
        timely decisions and respond proactively to dynamic conditions. By
        processing data as it is generated, airlines and air traffic controllers
        can anticipate issues, optimize operations, and improve safety.
        \item \textbf{Improve Operational Efficiency:} ncreasing operational
        efficiency is a primary goal, as it allows airlines to better manage
        their resources, streamline processes, and reduce delays. Enhanced
        efficiency leads to smoother operations, lower costs, and a more
        satisfactory experience for passengers.
        \item \textbf{Cost Optimization through Resource Management:} Effective
        resource management is crucial for cost optimization. By strategically
        allocating resources—such as fuel, crew, and equipment—airlines can
        minimize expenses without compromising service quality. Optimized costs
        contribute to a more sustainable business model and allow airlines to
        remain competitive in a challenging industry landscape.

    \end{itemize}
    $\quad$These objectives serve as the foundation for a transformative
    approach in the aviation industry, aiming to create a more responsive,
    efficient, and cost-effective operational framework powered by data-driven
    insights.
\subsection{Proposal}
To address the challenges in the aviation industry, this proposal focuses on
leveraging data from multiple sources, advanced analytics, and user-friendly
interfaces to improve decision-making, operational efficiency, and overall
service quality. The key components of the proposal include:
    \begin{itemize}
        \item \textbf{Real-time Data Ingestion from Multiple Sources} By
        integrating real-time data from various sources such as FlightAware,
        Flightradar24, and Flight Air Map, the aviation industry can access
        comprehensive and up-to-date information on flights, weather, and air
        traffic. This continuous data stream allows for improved situational
        awareness and facilitates faster responses to changing conditions.
        \item \textbf{Advanced Analytics and Predictive Models:} Employing
        advanced analytics and predictive modeling can help the industry move
        from reactive to proactive decision-making. Predictive models can
        forecast potential disruptions, optimize flight schedules, and support
        fuel and resource management. This data-driven approach enables better
        preparation for operational challenges, contributing to more efficient
        and reliable services.
        \item \textbf{User-friendly Dashboards for Reporting:} Developing
        intuitive and user-friendly dashboards allows stakeholders to easily
        interpret complex data. These dashboards can present key metrics,
        trends, and predictive insights, making it easier for decision-makers to
        act on the information at hand. Clear and accessible data visualization
        improves operational transparency and enhances collaboration across
        teams.
    \end{itemize}
    $\quad$This proposal aims to create a robust, data-centric system that
    supports real-time decision-making and predictive insights, ultimately
    leading to higher operational efficiency, cost savings, and improved
    customer satisfaction in the aviation industry.

\section{How the application works}
\subsection{Workflow}
The aviation data application operates through a structured workflow that
involves \textbf{Data Collection, Data Storage and Processing, and Analytics and
Reporting}. This approach ensures that vast amounts of real-time data are
collected, efficiently processed, and converted into actionable insights.

\subsection*{Step 1: Data collection}
The application’s workflow begins with \textbf{Data Collection}, where
information is gathered from multiple aviation-related sources such as
FlightAware, Flightradar24, Airline On-Time Statistics, and Flight Air Map.
These sources provide essential data, including flight schedules, delays, and
real-time positioning of aircraft. To create a comprehensive view of operations,
the application integrates this data in real-time from additional sources like
air traffic control systems and weather services, which helps in forming a
complete picture of flight conditions necessary for informed decision-making.

\subsection*{Step 2: Data storage and processing}
Once collected, data is organized and processed in the \textbf{Data Storage and
Processing} stage. Initially, raw data from diverse sources is stored in a data
lake (blob storage), which can accommodate various unstructured data types. To
enable real-time data transformation, the application employs PySpark, allowing
for the rapid processing of large datasets. The processed data is then stored in
a Yugabyte data warehouse, optimized for complex querying and analysis,
providing an efficient setup for generating analytical insights. This data is
structured using a star schema, a modeling technique that organizes information
into fact and dimension tables, which facilitates efficient retrieval and
supports detailed analysis.

\subsection*{Step 3: Analytics and reporting}
In the \textbf{Analytics and Reporting} stage, the application leverages machine
learning models to perform predictive analytics, such as forecasting maintenance
needs and potential flight delays by analyzing historical and real-time data.
This predictive capability allows the aviation industry to move from reactive to
proactive operations. Additionally, the application visualizes key performance
indicators (KPIs) and other metrics through user-friendly dashboards. These
visualizations make complex data accessible and interpretable, supporting quick,
data-driven decision-making and enhancing overall operational efficiency.

\subsection{Benefits to stakeholders}
The aviation data application provides significant benefits to various
stakeholder groups, each of which plays a crucial role in ensuring smooth and
efficient operations:

\begin{itemize}
    \item \textbf{Operational Teams:} With access to real-time insights,
    operational teams can improve scheduling and resource allocation. This
    capability allows for better alignment of flights, staff, and equipment,
    leading to enhanced efficiency and minimizing potential disruptions.
    \item \textbf{Maintenance Crews:} Management teams benefit from data-driven
    strategic planning based on both historical and real-time flight operation
    data. This comprehensive insight allows managers to make informed decisions,
    optimize resource allocation, and plan for the future, ultimately driving
    organizational success and resilience.
    \item \textbf{Management:} Management teams benefit from data-driven
    strategic planning based on both historical and real-time flight operation
    data. This comprehensive insight allows managers to make informed decisions,
    optimize resource allocation, and plan for the future, ultimately driving
    organizational success and resilience.
    \item \textbf{Passengers:} Passengers experience an enhanced travel
    experience due to fewer delays, optimized scheduling, and reduced
    cancellations. By improving operational efficiency and reliability, the
    application contributes to a smoother journey for travelers, increasing
    customer satisfaction and building loyalty.
\end{itemize}
$\quad$Overall, the application’s capabilities in real-time data processing,
predictive maintenance, and strategic planning empower each stakeholder group,
fostering a more efficient, proactive, and customer-centric aviation industry.

\section{Implementation plan}
The implementation plan for the aviation data application is divided into four
phases: \textbf{Planning and Requirements Gathering, Development, Testing, and
Deployment}. Each phase is designed to ensure that the application is built,
tested, and deployed effectively, meeting both technical and business
objectives.

\subsection*{Phase 1: Planning and Requirements Gathering}
This initial phase, lasting 1.5 weeks (Week 3 to mid-Week 4), involves defining
the project scope and objectives. It includes identifying the necessary data
sources, such as FlightAware, Flightradar24, and Flight Air Map, and outlining
infrastructure requirements. This foundational phase sets the direction for the
subsequent development work.

\subsection*{Phase 2: Development}
Spanning 2.5 weeks (mid-Week 4 to Week 7), this phase focuses on setting up the
big data infrastructure. This involves configuring Yugabyte for data
warehousing, establishing blob storage for the data lake, and utilizing PySpark
for distributed data processing. Additionally, data ingestion and processing
pipelines are developed to handle real-time data from various sources, ensuring
the system is equipped for efficient data handling.
\subsection*{Phase 3: Testing}
Lasting 2 weeks (Week 7 to Week 9), this phase includes pilot testing with
selected flight routes to validate data accuracy and system performance. Stress
testing is also conducted to ensure that the application can scale effectively
under heavy data loads, confirming its robustness before full deployment.
\subsection*{Phase 4: Deployment}
The final phase, taking 2 weeks (Week 9 to Week 11), involves the full-scale
implementation of the system, making it fully operational. This deployment phase
ensures that the application is accessible to stakeholders, with all
functionalities optimized for real-world use.

\section{Case studies and Examples}
    \subsection{Cebu Pacific’s Fuel Efficiency Program}
    Cebu Pacific, a major airline in the Philippines, implemented the
    \textbf{SkyBreathe} system — a big data and AI-driven solution focused on
    optimizing fuel consumption. This system leverages advanced analytics to
    track and analyze flight data, helping the airline make data-informed
    decisions that reduce fuel use. By fostering a data-driven culture and
    focusing on efficiency improvements, Cebu Pacific achieved significant cost
    savings.

    In 2017, Cebu Pacific saved over \textbf{11 million USD} by enhancing fuel
    efficiency through SkyBreathe. This achievement highlights the potential
    impact of big data and AI in the aviation industry, where even marginal
    improvements in fuel consumption can lead to substantial financial and
    environmental benefits. This case study serves as an example of how
    data-driven solutions can support sustainability goals and operational
    efficiency within the aviation sector.
    \subsection{Delta Airlines’ Use of Data Analytics}
    Delta Airlines successfully leveraged data analytics to enhance its baggage
    handling processes and improve customer engagement. By using data-driven
    insights, Delta was able to identify inefficiencies in baggage management
    and implement solutions that streamlined operations. This data-centric
    approach allowed the airline to track, predict, and reduce instances of
    mishandled baggage.

    As a result, Delta achieved a \textbf{71\% reduction in baggage
    mishandling}, significantly improving its operational efficiency and
    reducing customer frustration associated with lost or delayed luggage. This
    improvement also contributed to a noticeable increase in customer
    satisfaction scores, as passengers experienced smoother and more reliable
    travel experiences. Delta’s use of data analytics exemplifies how strategic
    data application can lead to tangible improvements in both operational
    performance and customer service in the aviation industry.
\section{Potential Challenges and Mitigation Strategies}
Implementing a data-driven application in the aviation industry comes with a range of challenges that must be addressed to ensure successful deployment and operation. Here are some of the key challenges and corresponding strategies for mitigating them:
\begin{itemize}
    \item \textbf{Data Security and Privacy:} Protecting sensitive information is critical in aviation, where breaches can have serious consequences. To mitigate these risks, the system should incorporate strong encryption methods, role-based access controls, and regular security audits. These measures help safeguard data and ensure that only authorized personnel can access sensitive information, reducing the risk of unauthorized access and data breaches.
    \item \textbf{Integration with Existing Systems:} Aviation companies often rely on a variety of legacy systems, making it challenging to integrate new solutions seamlessly. To address this, APIs and middleware solutions can be utilized to facilitate smooth communication between the new application and existing systems. This approach minimizes disruption and allows the new system to function harmoniously with the infrastructure already in place.
    \item \textbf{Change Management:} Transitioning to a data-driven system may face resistance from employees unfamiliar with new technologies or processes. Developing a comprehensive training and communication plan is essential for easing this transition. Providing training sessions, resources, and open channels of communication can help stakeholders understand the benefits of the new system, fostering acceptance and smoother adoption.
    \item \textbf{Regulatory Compliance:} The aviation industry is highly regulated, with strict standards around data governance and operational procedures. Ensuring adherence to aviation regulations and data governance standards is crucial. This can be achieved by continuously monitoring regulatory requirements and integrating compliance checks into the application’s operations. By aligning with industry standards, the application can avoid legal issues and maintain credibility.
\end{itemize}
$\quad$Addressing these challenges proactively allows for a smoother implementation process and ensures that the application operates effectively within the highly regulated and sensitive aviation environment.

\nocite{*}
\bibliographystyle{plainnat}

\bibliography{refs}

\end{document}